\section{Altitude Transitions in the Redux Framework}

\subsection{Overview}

The altitude system provides a structured hierarchy for organizing
geometric, structural, and analytical transformations within the Redux
pipeline. Each altitude corresponds to a distinct mathematical layer,
and transitions between altitudes represent the application of specific
operators that elevate the representation.

This section formalizes the altitude system and defines the transition
maps between altitudes.

\subsection{Altitude Definitions}

We define the following altitudes:

\begin{itemize}
    \item \textbf{A3}: Structural domain (degree normalization, symbolic input).
    \item \textbf{A5}: Spectral domain (eigenbasis representation).
    \item \textbf{A7}: Curvature domain (local geometric modulation).
    \item \textbf{A8}: Holonomy and topology domain (accumulated and global geometry).
    \item \textbf{A9}: Invariant domain (analytical diagnostics).
\end{itemize}

Each altitude corresponds to a distinct mathematical space and operator
family.

\subsection{Altitude Transition Maps}

Let $\mathcal{A}_k$ denote the space associated with altitude $k$. Each
transition is defined by a map
\[
T_{k \to k'} : \mathcal{A}_k \to \mathcal{A}_{k'}.
\]

The Redux pipeline uses the following transitions:

\begin{align*}
T_{3 \to 5} &: \text{Structural normalization} \\
T_{5 \to 7} &: \text{Spectral projection} \\
T_{7 \to 8} &: \text{Curvature modulation} \\
T_{8 \to 8'} &: \text{Holonomy transport} \\
T_{8' \to 8''} &: \text{Topology modulation} \\
T_{8'' \to 9} &: \text{Invariant evaluation}.
\end{align*}

These transitions form a directed chain:
\[
\mathcal{A}_3
    \xrightarrow{T_{3 \to 5}}
\mathcal{A}_5
    \xrightarrow{T_{5 \to 7}}
\mathcal{A}_7
    \xrightarrow{T_{7 \to 8}}
\mathcal{A}_8
    \xrightarrow{T_{8 \to 9}}
\mathcal{A}_9.
\]

\subsection{Compositional Structure}

The altitude transitions compose to form the full embedding map
\[
E = T_{8 \to 9}
    \circ T_{7 \to 8}
    \circ T_{5 \to 7}
    \circ T_{3 \to 5}.
\]

Thus, for any input structure $G$,
\[
E(G) \in \mathcal{A}_9.
\]

This compositional structure ensures that each altitude contributes a
distinct mathematical transformation.

\subsection{Monotonicity of Altitude}

Altitude transitions are monotonic: once a representation ascends to a
higher altitude, it does not return to a lower one. Formally,
\[
k < k' \implies
T_{k' \to k} \text{ is undefined}.
\]

This reflects the fact that each altitude adds information that cannot
be removed without loss.

\subsection{Altitude Stability}

Let $G$ and $G'$ be two input structures. The altitude system satisfies
the stability bound
\[
\| E(G) - E(G') \|
    \leq C \| G - G' \|,
\]
where $C$ depends on geometric bounds of the manifold. This ensures that
altitude transitions preserve continuity.

\subsection{Altitude Interpretation}

The altitude system provides a conceptual and mathematical framework for
understanding the Redux pipeline. Each altitude corresponds to a
distinct geometric or analytical layer, and transitions between altitudes
represent the application of operators that elevate the representation.

This section completes the formal definition of the altitude hierarchy.
